\chapter{修订与说明}
\label{cha:change}

本人在使用过程中发现了一些与《郑州大学学位论文写作规范》不一致的地方，因此根据原有的模板进行了一些修改。此外，根据本人的使用习惯增加和删除了部分宏包、注释等。最后给出了一些可能遇到的问题。由于能力有限，本版本依然有一些问题无法解决。

\rightline{——\color{teal}{$\mathsf{S} \mathsf{Q} \mathsf{K}$}   }

\section{编译环境的变化}
\subsection{一键盲审}
$\star$ 新增了盲审专用格式文件{\color{blue}MS.cls}。

$\star$ 在{\color{blue}main.tex}文件中修改环境即可

$\star$ 正常编译时为zzuthesis
\begin{latex}
    \documentclass[master]{zzuthesis}
\end{latex}

$\star$ 盲审编译时为MS
\begin{latex}
    \documentclass[master]{MS}
\end{latex}


\subsection{新定义命令}
$\star$ 文献引用命令，{\color{blue}\textbackslash mycite\{chen2001\}}编译后为\mycite{chen2001}，{\color{blue}\textbackslash cite\{chen2001\}}编译后为\cite{chen2001}。
\begin{latex}
    \newcommand{\mycite}[1]{\scalebox{1.3}[1.3]{\raisebox{-0.65ex}{\cite{#1}}}}
\end{latex}



\subsection{宏包的变化}

所有的宏包都在{\color{blue}{.cls}}格式文件中

$\star $ 支持中文字体——{\heiti 黑体}
\begin{latex}
    \RequirePackage{xeCJK}
\end{latex}

$\star$ 支持{\color{blue}公式}中字母字体为——Times New Roman ，并在{\color{blue}{main.tex}}文件中增加命令{\color{blue}\textbackslash setmainfont\{Times New Roman\}}
\begin{latex}
    \RequirePackage{mathptmx}
\end{latex}


$\star$ 支持{\color{blue}正文}中字母字体为——Times New Roman 
\begin{latex}
    \RequirePackage{fontspec}
\end{latex}

$\star$ 支持全文任意字体加粗
\begin{latex}
    \RequirePackage{bm}
\end{latex}

$\star$ 支持横向页面
\begin{latex}
    \RequirePackage{pdflscape}
\end{latex}

$\star$ 删除{\color{blue}\{unicode-math\}}宏包，使用{\color{blue}\{amsmath,amsthm\}}宏包。前者宏包会导致部分大写的希腊字母无法正体。
\begin{latex}
    \RequirePackage{amsmath,amsthm}
    %\RequirePackage{unicode-math}                        
    %\unimathsetup{                                       
    %  math-style = ISO,                                  
    %  bold-style = ISO,                                  
    %  nabla      = upright,                              
    %  partial    = upright,                              
    %} 
\end{latex}


\section{内容修订}
\subsection{封面}
\begin{itemize}
    \item 更换了新版的校徽
    \item 删除了作者信息部分的下划线，并保持左对齐
    \item 修改了{\color{blue}学位论文原创性声明}和{\color{blue}学位论文使用授权声明}的字体
\end{itemize}
\subsection{摘要}
\begin{itemize}
    \item 删除了关键词与摘要内容的空行
\end{itemize}
\subsection{目录}
\begin{itemize}
    \item 增加了中英文摘要、插图清单、表格清单和主要符号对照表的目录
    \item 各级标题与页码之间的点统一使用小四号
\end{itemize}
\subsection{正文}
\begin{itemize}
    \item 章节字体修改为——黑体
    \item 正文英文和公式的字体修改为——Times New Roman
\end{itemize}
\subsection{参考文献}
\begin{itemize}
    \item 卷号后的冒号修改为中文冒号“：”
    \item 所有逗号使用中文逗号“，”
    \item 网页引用的类型修改为——[EB]
\end{itemize}
\subsection{致谢}
\begin{itemize}
    \item 将致谢正文放置文章的最后
\end{itemize}

\section{Q \& A}

\noindent {\sihao\color{red} Q1: 由于图表中引用了文献导致参考文献无法正确排序}

修订者是使用{\color{blue}VS Code} + {\color{blue}texlive2020}进行编译
\begin{itemize}
    \item 每次编译前
    \item 需要先清理辅助文件(Clean up auxiliary files)
    \item 再进行编译(Recipe: xe->bib->xe->xe)
    \item 编译后不清理，以保证第二个问题的解决
\end{itemize}

Another Way:设置{\color{blue}VS Code}的{\color{blue}settings.json}文件，每次编译后默认删除{\color{blue}main.lof}  和 {\color{blue}main.lot} 文件。此方法在处理第二个问题时，需先关闭自动删除功能，将手动修改的两个文件保存至父目录，再打开功能。代码如下：
\begin{latex}
    "latex-workshop.latex.autoClean.run": "onBuilt", //注意结尾是 t 不是 d
    "latex-workshop.latex.clean.fileTypes": [        //注意，'//'代表注释符
//  "*.lof",                           //此处取消注释符即开启自动删除功能 
//  "*.lot",                           //此处取消注释符即开启自动删除功能
//  "*.aux",
//  "*.bbl",
//  "*.blg",
//  "*.idx",
//  "*.ind",
//  "*.out",
//  "*.toc",
//  "*.acn",
//  "*.acr",
//  "*.alg",
//  "*.glg",
//  "*.glo",
//  "*.gls",
//  "*.ist",
//  "*.fls",
//  "*.log",
//  "*.fdb_latexmk",
//],
\end{latex}

如果使用其他软件+{\color{blue}texlive}编译，如{\color{blue}WinEdt}
\begin{itemize}
    \item 每次编译前
    \item 需要先清理辅助文件(点击垃圾箱图标)
    \item 再进行编译(XeTex->BibTex->XeTex->XeTex)
    \item 编译后不清理，以保证第二个问题的解决
\end{itemize}

没有删除辅助文件快捷键的软件，可手动进入{\color{blue}main.tex}文件目录下删去除{\color{blue}main.tex}文件外所有前缀为{\color{blue}main.}的文件。

\noindent {\sihao\color{red} Q2: 图表目录生成时会自动使用`caption'内全部内容  }
\begin{itemize}
    \item 为保证目录的简洁                            
    \item 可分别在`main.lof'和`main.lot'文件内手动修改
    \item 修改后需将文件保存至{\color{blue}main.tex}文件的父目录                  
    \item 每次通过Q1中的编译方法编译后
    \item 再将两个文件移入{\color{blue}main.tex}文件目录下
    \item 再次XeTex编译即可
\end{itemize}
